% ============================================================
%  ch03  第一讲 §1.4 Beauville--Laszlo 定理、§1.5 行列式线丛
%  原文行 1033–1191
% ============================================================

\subsection{Beauville--Laszlo 定理}

$\Gr_{\underline G}$ 还有一个重要的模空间解释，归功于 Beauville--Laszlo。

设 $X$ 是 $k$ 上既约连通曲线（不必光滑），$x\in |X|$ 是光滑闭点，$\underline G$ 是 $X$ 上光滑仿射群概形。定义预层
\begin{equation}\label{aff glob}
\Gr_{\underline G,x}(R)=\left\{(\mE,\beta)\ \left|\ \begin{array}{l}
\mE \text{ 是 } X_R \text{ 上的 } \underline G\text{-主丛},\\
\beta:\mE|_{X_R^*}\cong \mE^0|_{X_R^*} \text{ 是一个平凡化}
\end{array}\right.\right\}.
\end{equation}
令 $\underline G_x=\underline G\otimes_X \mO_x$。把 $X_R$ 上的 $\underline G$-主丛限制到 $D_{x,R}$，给出自然态射
\[
\on{res}:\Gr_{\underline G,x}\to \Gr_{\underline G_x}.
\]

\begin{thm}\label{glob aff}
上述态射是同构。
\end{thm}

同样，定理的证明有技术性，但想法很简单，先作非正式说明。若 $\mE$ 是 $D_x$ 上带平凡化 $\beta$（在 $D_x^*$ 上）的 $\underline G_x$-主丛，则可用 $\beta$ 把 $\mE$ 与 $X^*$ 上平凡的 $\underline G|_{X^*}$-主丛 $\mE^0$ 粘起来。这一构造给出 $\on{res}$ 的逆映射。

要把这个想法写成严格证明，先要讨论 Beauville--Laszlo 的下降定理（见 \cite{BL2}）。

\begin{thm}\label{descent lemma}
设 $p: A\to \tilde{A}$ 是交换环同态，$f\in A$。设 $f$（相应地，$p(f)$）在 $A$（相应地，$\tilde{A}$）中不是零因子，且诱导映射 $A/f\to \tilde{A}/p(f)$ 是同构。以 $\mM_f(A)$ 记无 $f$-挠的 $A$-模范畴，以 $\mM_{p(f)}(\tilde{A})$ 记 $(\tilde{A},p(f))$ 的类似范畴。以 $\mM$ 记三元组 $(M_1,M_2,\varphi)$ 的范畴，其中 $M_1$ 是 $A_f$-模，$M_2\in \mM_{p(f)}(\tilde{A})$，而 $\varphi$ 是 $M_1\otimes_{A_f}\tilde{A}_{p(f)}$ 与 $M_2\otimes_{\tilde{A}}\tilde{A}_{p(f)}$ 之间的同构。则
\begin{enumerate}
\item 函子 $M\mapsto M\otimes_A\tilde{A}$ 把 $\mM_f(A)$ 映入 $\mM_{p(f)}(\tilde{A})$，故有函子 $F : \mM_f(A)\to \mM$，它把 $M\in\mM_f(A)$ 送到 $(M_f, M\otimes_A\tilde{A}, \varphi)$，其中 $\varphi$ 是 $M_f\otimes_{A_f}\tilde{A}_{p(f)}$ 与 $(M\otimes_A\tilde{A})_{p(f)}$ 之间的自然同构。
\item $F : \mM_f(A)\to \mM$ 是范畴等价。
\item $M\in\mM_f(A)$ 是有限（相应地，平坦、有限投射）$A$-模，当且仅当 $M_f$ 与 $M\otimes_A\tilde{A}$ 都具有该性质。
\end{enumerate}
\end{thm}

在 \cite{BL2} 中取 $\tilde{A}=\hat A$，即 $A$ 关于 $f$ 的完备化；但同样的论证在此也适用。

\begin{rmk}
这个定理的内容是从覆盖 $(A_f, \tilde A)$ 把模下降到 $A$。它不能由通常的忠实平坦下降推出，因为没有假定 $A\to \tilde A$ 平坦。例如 $A$ 非 Noether 时，$A\to \hat A$ 通常不平坦。但即使在 Noether 情形，同构 $\varphi$ 也没有给出完整的下降数据，即不含 $\hat R\otimes_R\hat R$ 上的同构。
\end{rmk}

下面是该定理的推论。

\begin{cor}\label{global interpretation of loop}
环路群 $L\underline G_x$ 表示下述函子：它给每个 $k$-代数 $R$ 配上三元组 $(\mE,\al,\beta)$ 的集合，其中 $\mE$ 是 $X_R$ 上的 $\underline G$-主丛，$\al$ 是 $\mE|_{D_{x,R}}$ 的平凡化，$\beta$ 是 $\mE|_{X^*_R}$ 的平凡化。
\end{cor}

\begin{proof}
详细论证见 \cite[Lemma 6.1]{PZ}，思路如下。

可设 $X=\Spec A$ 仿射，$x$ 由 $t=0$ 定义。设 $R$ 是 $k$-代数，$g\in L\underline G_x(R)$。以 $H_1$（相应地，$H_2$）记 $\underline G\times_X D_{x,R}$（相应地，$\underline G\times_XX^*_R$）的正则函数环，它是平坦 $R[[t]]$-代数（相应地，$(A[t^{-1}]\otimes R)$-代数）。元素 $g$ 诱导代数同构 $\phi:H_1\otimes_{R[[t]]} R((t))\cong H_2 \otimes_{(A[t^{-1}]\otimes R)} R((t))$。

把定理 \ref{descent lemma} 用于三元组 $(H_1,H_2,\phi)$，得到平坦的 $(A\otimes R)$-代数 $H$。于是 $\mE=\Spec H$ 就是 $X_R$ 上的 $\underline G$-主丛，并带有典范平凡化 $\al$ 与 $\beta$。

反之，给定 $(\mE,\al,\beta)$，由 $g=(\beta|_{D_{x,R}^*})(\al|_{D_{x,R}^*})^{-1}$ 得到 $g$。
\end{proof}

现在比较命题 \ref{I:quotient} 与推论 \ref{global interpretation of loop}，定理 \ref{glob aff} 即得。也可以仿照上面的论证，直接构造 $\on{res}$ 的逆映射。细节留给读者。

\begin{rmk}
回忆 $\Aff_k$ 上 fpqc 层范畴与 $\on{Sch}_k$ 上 fpqc 层范畴等价。在此等价下，\eqref{aff glob} 中 $\Gr_{\underline G,x}$ 的模空间描述可自然推广到所有 $k$-概形。反之，把 \eqref{loc aff} 的模空间描述推广到所有 $k$-概形则很别扭：对非仿射概形 $S$，$S\hat{\times} \Spec k((t))$ 并不真的有意义。
\end{rmk}

\subsection{行列式线丛}\label{determinant line}

下一个主题是在 $\Gr_{\underline G}$ 上构造一个极丰沛线丛。回忆定理 \ref{rep} 证明中的交换图
\[\begin{CD}
\Gr^{(N)}@>>>\Gr^{(N+1)}\\
@VVV@VVV\\
\Gr(2nN)@>>>\Gr(2n(N+1))
\end{CD}\]
在通常的 Grassmann 簇 $\Gr(m)$ 上有行列式线丛 $\mL$：它把子空间 $L\subset k^m$ 配上商空间 $k^m/L$ 的最高外幂。这实际上是 $\Gr(m)$（各连通分支的）Picard 群的丰沛生成元。这些线丛在嵌入 $\Gr(m)\subset \Gr(m+1)$ 下相容，因此在 $\Gr_{\GL_n}$ 上定义了丰沛线丛 $\mL_{\det}$。我们更愿意给出一个更典范的构造。

先引入一个重要概念。

\begin{dfn}
若一个拓扑向量空间是某个离散向量空间的拓扑对偶（即通常的对偶，其 $0$ 的开邻域基由有限维子空间的零化子组成），则称它\textbf{线性紧}。若一个拓扑空间有一组由线性紧子空间组成的 $0$ 的邻域基，则称它\textbf{线性局部紧}，或\textbf{Tate 空间}。Tate 向量空间 $V$ 中的\textbf{格}是指 $V$ 的线性紧开子空间。
\end{dfn}

Tate 向量空间中任意两个格 $L_1, L_2$ 是\textbf{可公度的}，即商 $L_1/(L_1 \cap L_2)$ 与 $L_2/(L_1 \cap L_2)$ 都有限维。

\begin{rmk}\label{exact str}
Tate 向量空间构成范畴，其 Hom 为连续线性映射。这是 Quillen 意义下的正合范畴：容许单射为闭嵌入，容许满射为开满射。

Tate 向量空间还有不借助拓扑的范畴化途径。回忆对正合范畴 $\mC$，Beilinson 构造了「局部紧对象」范畴 $\underleftrightarrow{\lim} \mC$，作为 $\mC$ 的 pro-ind-对象范畴的满子范畴（\cite[Appendix]{Be}）。若把 $k$ 上有限维向量空间范畴 $\on{Vect}_k$ 看作正合范畴，则 Tate 向量空间范畴恰是 Beilinson 范畴 $\underleftrightarrow{\lim}\on{Vect}_k$。这一途径的好处是不用拓扑，并可推广到其他情形。
\end{rmk}

设 $V$ 是 $k$ 上 Tate 向量空间。对 $k$-代数 $R$，定义拓扑 $R$-模
\[V\hat\otimes_k R= \varprojlim_L (V/L)\otimes_kR,\]
其中反向极限取遍 $V$ 中的格。于是可定义群预层 $\GL(V)$，其 $R$-点为 $V\hat\otimes_k R$ 的连续 $R$-线性自同构。

\begin{dfn}
$V$ 中的一个 \textbf{$R$-族格}是指开有界子模 $\Lambda\subset V\hat\otimes_k R$，且 $V\hat\otimes_k R/\Lambda$ 是投射 $R$-模。
\end{dfn}

\begin{dfn}
\textbf{Sato Grassmann 簇} $\Gr(V)$ 是这样的预层：它给每个 $k$-代数 $R$ 配上 $V$ 中全体 $R$-族格构成的集合。
\end{dfn}

与定理 \ref{rep} 完全类似，
\begin{thm}
Sato Grassmann 簇 $\Gr(V)$ 可由归纳射影概形表示。
\end{thm}
事实上，用类似（而且更简单）的论证有自然同构
\[\Gr(V)\cong \varinjlim_{L_1\subset L_2}\Gr(L_2/L_1),\]
其中 $L_1\subset L_2$ 是两个格，故 $L_2/L_1$ 是有限维 $k$-向量空间，而 $\Gr(L_2/L_1)$ 表示分类 $L_2/L_1$ 的有限维子空间的通常 Grassmann 簇。

下面在 $\Gr(V)\times \Gr(V)$ 上引入行列式线丛（空间上线丛的定义见 \S\ref{Intro:sp}）。先回忆：有限投射 $R$-模 $M$ 的最高外幂记作 $\det(M)$，它是可逆 $R$-模。对 $\Lambda_1,\La_2\in\Gr(V)(R)$，定义 $(\La_1,\La_2)$ 的\textbf{相对行列式线}为
\[\det(\La_1|\La_2)=\det(L\hat\otimes R/\La_1)\otimes \det(L\hat\otimes R/\La_2)^{-1},\]
其中 $L$ 是 $V$ 中某个格，使得 $(L\hat\otimes R)/\La_1$ 与 $(L\hat\otimes R)/\La_2$ 都是投射 $R$-模。这在典范同构意义下与 $L$ 的选取无关。让 $R$ 变动，得到 $\Gr(V)\times\Gr(V)$ 上的线丛 $\mL_{\det}$。下面这个引理虽然初等，却概括了相对行列式线的重要性质。

\begin{lem}~
\begin{enumerate}
\item 对任意 $g\in  \GL(V)(R)$ 与 $\La_1,\La_2\in\Gr(V)(R)$，有典范同构 $\det(g\La_1|g\La_2)\cong\det(\La_1|\La_2)$，且对 $g,g'$，同构 $\det(gg'\La_1|gg'\La_2)\cong\det(g'\La_1|g'\La_2)\cong\det(\La_1|\La_2)$ 与 $\det(gg'\La_1|gg'\La_2)\cong\det(\La_1|\La_2)$ 一致。换言之，$\GL(V)$ 在 $\Gr(V)\times\Gr(V)$ 上的对角作用提升为 $\mL_{\det}$ 上的作用。
\item 对任意 $\La_1,\La_2,\La_3$，有典范同构 $\ga_{123}:\det(\La_1|\La_2)\otimes\det(\La_2|\La_3)\cong\det(\La_1|\La_3)$，使得对任意 $\La_1,\La_2,\La_3,\La_4$ 有 $\ga_{134}\ga_{123}=\ga_{124}\ga_{234}$。换言之，若以 $\mL_{\det}^\times$ 记 $\mL_{\det}$ 的全空间，则得到群胚
\[\mL_{\det}^\times\rightrightarrows\Gr(V).\]
\end{enumerate}
\end{lem}

\begin{rmk}
该群胚定义了 $\bG_m$-gerbe
\[\mD_V=[\Gr(V)/ \mL_{\det}^\times],\]
通常称为 $V$ 的\textbf{行列式 gerbe}（这个术语由 Kapranov 首先引入）。
\end{rmk}

\begin{rmk}
上述构造可以族化。确切地说，对交换环 $R$，有 Tate $R$-模范畴 $\on{Tate}_R$，它定义为 $\underleftrightarrow{\lim}\on{Proj}_R$ 的幂等完备化，其中 $\on{Proj}_R$ 是有限投射 $R$-模范畴（拓扑途径见 \cite{Dr}）。对 Tate $R$-模 $M$，有相应的 Sato Grassmann 簇 $\Gr(M)\to \Spec R$（见\emph{同上}），它可由归纳代数空间表示。此外，$\Spec R$ 上有行列式 gerbe $\mD_M$。
\end{rmk}

\begin{rmk}\label{graded line}
事实上，上述构造应当升级为分次版本。确切地说，对有限投射 $R$-模，可以把 $\det(M)$ 定义为分次可逆 $R$-模，其中 $\det(M)$ 的次数取为 $M$ 的秩（它是 $\Spec R$ 上局部常值函数）。此时上面定义的 $\det(\La_1\mid\La_2)$ 成为 $\Spec R$ 上的分次线丛。记住分次对行列式线上的分解结构至关重要（见注 \ref{graded factorization} 与注 \ref{det factorization2}），而后者对符号与互反律的应用又很重要（见 \cite{BBE, OZ1, OZ2}）。
\end{rmk}

取 $V=k((t))^n$，配以通常的 $t$-进拓扑。它是 Tate 向量空间，含格 $\Lambda_0=k[[t]]^n$。由引理 \ref{projectivity} 得典范闭嵌入 $i:\Gr_{\GL_{n}}\to\Gr(V)$。于是通过拉回得到线丛 $\mL_{\det}|_{\{\La_0\}\times\Gr_{\GL_n}}$，仍记作 $\mL_{\det}$。注意这是 $\Gr_{\GL_n}$ 上的极丰沛线丛。

一般地，设 $\rho: \underline G\to \GL_n$ 是命题 \ref{lin rep} 中忠实的 $\mO$-线性表示。则由命题 \ref{gen G}，它诱导局部闭嵌入 $f_\rho: \Gr_{\underline G}\to \Gr_{\GL_n}$，而 $\mL_{\det}$ 的拉回定义了 $\Gr_{\underline G}$ 上的极丰沛线丛。

\begin{rmk}\label{Pfaffians}
设 $\on{char} k\neq 2$（去掉这一假设见 \cite[\S\ 4.2.16]{BD}）。设 $V$ 是 $k$ 上 Tate 向量空间。$V$ 上的连续对称双线性型 $B$ 称为\textbf{非退化}的，若诱导映射 $V\to V^*$ 是同构，其中 $V^*$ 是 $V$ 的拓扑对偶。此时对格 $\Lambda$ 可定义其正交补 $\Lambda^\perp$。若 $\Lambda=\Lambda^\perp$，则称 $\Lambda$ 为\textbf{拉格朗日格}。正交 Sato Grassmann 簇 $\OGr(V,B)\subset \Gr(V)$ 是分类拉格朗日格的闭子空间。如 \cite[\S\ 4]{BD} 所构造，行列式线丛 $\mL_{\det}$ 在 $\OGr(V,B)\times \OGr(V,B)$ 上的限制有典范平方根 $\mL_{\on{Pf}}$，称为 \textbf{Pfaffian 线丛}。
\end{rmk}
